% META: {"id": "1NSI-PM-EX-004", "chapitre": "1NSI-PROJET-METHODES", "type_objet": "exercice", "capacites": ["BO-PREAMBULE-COMPETENCES-ORALES"], "parcours": 2, "competences": ["communiquer", "analyser"], "duree_min": 10, "mode_creation": "ex_nihilo", "sources_inspiration": [], "fichier_tex": "chapitres/1NSI-PROJET-METHODES/exercices/1NSI-PM-EX-004.tex", "corrige_tex": "chapitres/1NSI-PROJET-METHODES/corriges/1NSI-PM-CO-004.tex", "status": "generated"}
% BEGIN-VERIFY
% def a_une_docstring(fonction):
%     return fonction.__doc__ is not None and len(fonction.__doc__.strip()) > 0
% def tri_insertion(tableau):
%     tableau = tableau[:]
%     for i in range(1, len(tableau)):
%         valeur = tableau[i]
%         j = i - 1
%         while j >= 0 and tableau[j] > valeur:
%             tableau[j + 1] = tableau[j]
%             j -= 1
%         tableau[j + 1] = valeur
%     return tableau
% assert a_une_docstring(tri_insertion) is False
% def tri_insertion_documente(tableau):
%     """Trie tableau par ordre croissant, par insertion.
%
%     Precondition : tableau est une liste d'elements comparables entre eux.
%     Postcondition : le resultat est une liste triee, de meme longueur.
%     """
%     tableau = tableau[:]
%     for i in range(1, len(tableau)):
%         valeur = tableau[i]
%         j = i - 1
%         while j >= 0 and tableau[j] > valeur:
%             tableau[j + 1] = tableau[j]
%             j -= 1
%         tableau[j + 1] = valeur
%     return tableau
% assert a_une_docstring(tri_insertion_documente) is True
% assert tri_insertion_documente([3, 1, 2]) == [1, 2, 3]
% END-VERIFY
\begin{exercice}{1NSI-PM-EX-004}{2}{10}
On reprend la fonction \lstinline{tri_insertion} d'un chapitre précédent, sans aucune
documentation :
\begin{python}
def tri_insertion(tableau):
    tableau = tableau[:]
    for i in range(1, len(tableau)):
        valeur = tableau[i]
        j = i - 1
        while j >= 0 and tableau[j] > valeur:
            tableau[j + 1] = tableau[j]
            j -= 1
        tableau[j + 1] = valeur
    return tableau
\end{python}
\begin{enumerate}
  \item Vérifier, avec \lstinline{a_une_docstring} du cours, que cette fonction n'est pas documentée.
  \item Écrire une docstring complète pour cette fonction (rôle, précondition, postcondition), puis vérifier qu'elle est bien détectée comme documentée.
  \item Si tu devais présenter cette fonction en $2$ minutes à l'oral, quels seraient les $3$ points les plus importants à expliquer (au-delà d'une lecture ligne par ligne du code) ?
\end{enumerate}
\end{exercice}
